\documentclass{article}


\usepackage[utf8]{inputenc}
\usepackage[spanish]{babel} 
\usepackage{graphicx}
\usepackage{amssymb}

\title{Explicación en lenguaje natural}
\author{}
\date{}
\setlength{\parskip}{10pt}


\begin{document}
\maketitle

\section{Introducción}

\subsection{Motivación}

La mayoría de estudiantes enfrentan dificultades para gestionar adecuadamente su hogar debido a la falta de tiempo, ya que gran parte de su jornada
está dedicada a sus estudios. Esta situación convierte la organización y la realización de tareas domésticas en un problema, pues estas actividades requieren 
una planificación cuidadosa y tiempo para llevarlas a cabo. Por otro lado, la gestión de los gastos compartidos puede resultar especialmente compleja en un entorno 
de estudiantes, ya que es común que cada miembro del hogar se encargue de diferentes pagos. Esto puede generar confusiones y equivocaciones, 
respectoa a quién ha pagado ciertos gastos y qué cantidades están pedientes, lo que agrava aún más la dificultad de una correcta
administación del hogar.

Después de 5 años viviendo de manera independiente, he experimentado de primera mano las dificultades de gestionar un hogar compartido con otros estudiantes o trabajadores. En numerosas ocasiones 
el tiempo libre del que disponíamos no se ha destinado a la organización del hogar. Aunque las tareas se asignaban semanalmente, no siempre se cumplían, ya sea
por falta de tiempo y o por descuidos. Además, la administación de los gastos ha sido, en muchas ocasiones, caótica, ya que no siempre se llevaba un control adecuado sobre quién ha pagado qué,
y algunas personas solían olvidar devolver el dinero a tiempo.  

Actualmente no existe ninguna aplicación que ofrezca una solución integral tanto para la gestión de las tareas domesticas como para la gestión de los gastos. Las aplicaciones
disponibles abordan estos aspectos por separado. Además la mayoría de ellas están orientadas a hogares familiares donde la distribución de las tareas no se gestiona de la misma forma que 
en los pisos de estudiantes. Por eso es fundamental desarrollar una aplicación que integre todas estas funcionalidades, dirigida específicamente a estudiantes o trabajadores que compartan vivienda. 

\subsection{Objetivos del proyecto}

El objetivo general de este proyecto es facilitar a los estudiantes la organización y gestión eficiente de su hogar, ofreciendo soluciones que aborden las dificultades 
que puedan surgir en la planificación y distribución de responsabilidades. Para poder alcanzar este objetivo, podemos diferenciar una serie de objetivos específicos:


\subsection{Estructura del documento}



\subsection{Problema a resolver}

Como se mencionó en el apartado anterior, el principal problema que esta aplicación 
busca resolver es la falta de organización con la gestión del hogar. No obstante,
también pretende abordar los problemas derivados de una mala gestión de dichas tareas, 
como la falta de comunicación entre los miembros del hogar, el descontrol en la administración de los gastos, 
la falta de tiempo para llevar a cabo las actividades domésticas, o el desordeb que surge cuando alguna de estas tareas 
no se completan.




\section{Descripción General}

\subsection{Funcionalidad Principal}
\begin{enumerate}
    \item \textbf{Lista de tareas domesticas compartida}. La aplicación permitirá a los miembros del hogar definir una lista de tareas domésticas que se deben realizar semanalmente. Cada tarea la tendrá que 
    realizar un miembro del hogar distinto, y se podrá asignar de manera aleatoria o manual. Además, se podrá establecer una fecha límite para la realización de cada tarea.
    \item \textbf{Gestión de gastos compartidos}. Contará con una funcionalidad que permita a todos los miembros del hogar registrar los gastos que han realizado y a quién se le debe dinero.
    \item \textbf{Recordatorios Automáticos}. La aplicación enviará recordatorios automáticos a los miembros del hogar para que realicen las tareas domésticas asignadas o para que paguen los gastos pendientes.
    \item \textbf{Lista de la compra compartida}. Se podrá crear una lista de la compra compartida en la que todos los miembros del hogar podrán añadir los productos que necesitan.
    \item \textbf{Calendario de tareas}. La aplicación contará con un calendario en el que se mostrarán las tareas domésticas asignadas a cada miembro del hogar y las fechas límite para su realización.
\end{enumerate}

\subsection{Público Objetivo}
Mi aplicación está dirigida a estudiantes o trabajadores que compartan un hogar con otras personas y que necesiten una herramienta para organizar y gestionar las tareas domésticas y los gastos compartidos.
Este grupo se beneficiaría especialmente de la aplicación, ya que la falta de tiempo y la necesidad de coordinación son problemas comunes en este tipo de hogares.

\section{Impacto Esperado}

Espero que esta aplicación permitia a los estudiantes ahorrar tiempo en la organización y asignación de las tareas domésticas, 
evitando además la necesidad de estar constantemente verificando si cada miembro ha cumplido con sus responsabilidades. 
Asimismo, espero que la aplicación contribuya a eliminar los problemas relacionados con retrasos en los pagos en muchos hogares compartidos, 
y que facilite un mayor control sobre quién ha efectuado cada pago.


Aunque la aplicación ha sido diseñada específicamente para estudiantes, su utilidad podría extenderse a otros grupos o tipos de hogares, 
como trabajadores que compartan vivienda o familias. 

\section{Conclusión}
 
En conclusión, esta aplicación pretende ser una herramienta útil para la organización y gestión de tareas domésticas 
y gastos compartidos en hogares de estudiantes. Espero que la aplicación contribuya a mejorar la convivencia en estos hogares, 
facilitando la comunicación y la coordinación entre los miembros del hogar, y permitiendo un mayor control sobre las tareas y los gastos.






\end{document}